\documentclass[11pt]{article}
\newcommand{\paso}{3}
\usepackage[T1]{fontenc}
\usepackage[utf8]{inputenc}
\usepackage[spanish,es-nodecimaldot,es-noquoting,es-lcroman]{babel}
\spanishplainpercent
\usepackage{lmodern}
\usepackage[a4paper,margin=21mm,headheight=15pt]{geometry}
\usepackage{amsmath,amssymb,bm}
\usepackage{booktabs,array}
\usepackage{tikz}
\usetikzlibrary{arrows.meta,positioning,calc}
\usepackage[most]{tcolorbox}
\usepackage{microtype}
\usepackage{fancyhdr}
\usepackage{hyperref}
\hypersetup{colorlinks=true,linkcolor=blue!45!black,urlcolor=blue!45!black}
\setlength{\parindent}{0pt}
\setlength{\parskip}{5pt}
\setlength{\emergencystretch}{2em}
\pagestyle{fancy}
\fancyhf{}
\fancyhead[L]{\small BEM del cilindro: paso \paso}
\fancyhead[R]{\small Guía matemática independiente}
\fancyfoot[C]{\thepage}
\newcommand{\vct}[1]{\bm{#1}}
\newcommand{\dd}{\,\mathrm{d}}
\newcommand{\ii}{\mathrm{i}}
\newcommand{\norm}[1]{\left\lVert#1\right\rVert}
\newcommand{\VP}{\operatorname{VP}}
\newcommand{\tr}{\operatorname{tr}}
\newtcolorbox{idea}[1]{breakable,colback=blue!3,colframe=blue!50!black,
  title={#1},fonttitle=\bfseries,boxrule=0.7pt}
\newtcolorbox{alerta}[1]{breakable,colback=orange!4,colframe=orange!65!black,
  title={#1},fonttitle=\bfseries,boxrule=0.7pt}
\newtcolorbox{resultado}[1]{breakable,colback=teal!3,colframe=teal!55!black,
  title={#1},fonttitle=\bfseries,boxrule=0.7pt}
\newcommand{\inicio}[2]{%
 {\small\color{blue!55!black}\bfseries SERIE: LOS SIETE PASOS MATEMÁTICOS DEL BEM}\par
 {\LARGE\bfseries Paso \paso. #1\par}\vspace{3mm}
 \begin{idea}{¿Qué busca este paso?}#2\end{idea}}
\newcommand{\mapa}[3]{%
 \begin{center}
 \begin{tikzpicture}[node distance=4mm,
   caja/.style={draw=blue!50!black,fill=blue!3,rounded corners=2pt,
     text width=13cm,align=center,inner sep=6pt},
   >={Stealth[length=2mm]}]
 \node[caja] (a){#1};
 \node[caja,below=of a] (b){#2};
 \node[caja,below=of b] (c){#3};
 \draw[->,thick,blue!60!black] (a)--(b);
 \draw[->,thick,blue!60!black] (b)--(c);
 \end{tikzpicture}
 \end{center}}
\newcommand{\base}{%
 \smallskip{\footnotesize\color{black!65}
 Modelo de referencia: cilindro sólido, material homogéneo e isótropo,
 deformación plana, pequeñas deformaciones, sin fuerzas de volumen;
 amplitudes complejas con \(e^{+\ii\omega t}\), \(\omega>0\).}\par}
\newcommand{\fuentes}{%
 \par\medskip\begingroup\footnotesize
 \textbf{Referencia y continuación.}
 Este documento desarrolla una etapa del esquema
 \texttt{esquema\_ecuaciones\_bem.tex}. Las ecuaciones y los ejemplos se
 contrastaron con \texttt{bem\_cilindro.py}, ambos en la carpeta superior.
 La guía general y sus deducciones ampliadas se encuentran en
 \texttt{../guia\_bem/guia\_bem.pdf}. No se modifica la implementación.
 \par Para recompilar este documento, mantener \texttt{estilo.tex} en
 la misma carpeta y ejecutar \texttt{pdflatex} sobre su archivo
 \texttt{.tex}, desde esa carpeta.\par\endgroup}

\begin{document}
\inicio{La identidad de Green}{Reemplazar la ecuación diferencial en el área
por una relación entre desplazamientos y tracciones de su contorno.
Este es el cambio conceptual central del BEM: no aproximar la ecuación en
cada punto interior, sino aprovechar una solución fundamental conocida.}
\base

\mapa{Campo real: \(\nabla\!\cdot\boldsymbol\sigma+\rho\omega^2\vct u=\vct0\)}
{Compararlo con un campo auxiliar generado por una fuerza puntual}
{Somigliana: \(\vct u(\vct\xi)=\int_\Gamma U\vct t\,\dd\Gamma
 -\int_\Gamma T\vct u\,\dd\Gamma\), si \(\vct\xi\in\Omega\)}

\section*{1. Las dos posiciones no son intercambiables}
\begin{center}
\begin{tikzpicture}[>={Stealth[length=2mm]}]
 \fill[blue!4] (0,0) ellipse(2.4 and 1.4);
 \draw[thick,blue!55!black] (0,0) ellipse(2.4 and 1.4);
 \fill (-0.8,0) circle(2pt) node[below]{$\vct\xi$: fuente fija};
 \fill (1.7,0.99) circle(2pt) node[above left]{$\vct x\in\Gamma$};
 \draw[->,thick] (-0.8,0)--(1.7,0.99) node[midway,above]{$\vct x-\vct\xi$};
 \draw[->,teal!60!black,thick] (1.7,0.99)--(2.3,1.7) node[right]{$\vct n(\vct x)$};
 \node at (0.9,-0.6){$\Omega$};
 \node[text width=6cm,align=left] at (6,0){
   Se fija $\vct\xi$.\\
   $\vct x$ recorre el contorno.\\
   La normal se evalúa en $\vct x$,\\
   no en la fuente interior.};
\end{tikzpicture}
\end{center}

\section*{2. Qué representan \(U\) y \(T\)}
Para una fuerza auxiliar unitaria en dirección \(a\), aplicada en
\(\vct\xi\):
\[
 v_j(\vct x)=U_{aj}(\vct x,\vct\xi),\qquad
 t_j^{\mathrm{aux}}(\vct x)=T_{aj}(\vct x,\vct\xi).
\]
En estos documentos el \textbf{primer índice identifica la fuerza auxiliar}
y el segundo la componente del campo. Así, la ecuación para \(u_a\)
contiene \(U_{aj}t_j\) y \(T_{aj}u_j\).
Es la convención que reproduce los bloques de la implementación.

\begin{alerta}{Aclaración de índices respecto al esquema resumido}
Para obtener \(T_{aj}\), se aplica el operador de tracción a la
\emph{fila} \(U_{a\cdot}\). Usar en cambio la columna y conservar los mismos
índices produce la matriz traspuesta. La simetría de \(U\) puede ocultar
esta diferencia; \(T\) no es simétrica en general.
\end{alerta}

\newpage
\section*{3. Comparar dos estados del mismo material}
Sean el estado físico \((\vct u,\boldsymbol\sigma)\) y el estado auxiliar
\((\vct v,\boldsymbol\tau)\), a la misma frecuencia. Ambos comparten el
tensor elástico complejo \(C_{ijkl}\), con simetría mayor
\(C_{ijkl}=C_{klij}\). Por esa simetría y la de los esfuerzos:
\[
 u_{i,j}\tau_{ij}
 =\varepsilon_{ij}(\vct u)\tau_{ij}
 =\varepsilon_{ij}(\vct u)C_{ijkl}\varepsilon_{kl}(\vct v)
 =v_{i,j}\sigma_{ij}.
\]
Esta igualdad es el núcleo de la reciprocidad.

\subsection*{3.1 La regla del producto crea una divergencia}
\[
 \begin{aligned}
 \partial_j(u_i\tau_{ij}-v_i\sigma_{ij})
 &=u_{i,j}\tau_{ij}-v_{i,j}\sigma_{ij}\\
 &\quad+u_i\tau_{ij,j}-v_i\sigma_{ij,j}\\
 &=u_i\tau_{ij,j}-v_i\sigma_{ij,j}.
 \end{aligned}
\]
El primer par de términos se cancela por reciprocidad. Para el campo real
sin fuerzas de volumen:
\[
 \sigma_{ij,j}=-\rho\omega^2u_i.
\]
Para el auxiliar, con fuerza puntual unitaria en dirección \(a\):
\[
 \tau_{ij,j}=-\rho\omega^2v_i-\delta_{ia}\delta(\vct x-\vct\xi).
\]
La delta de Dirac concentra una fuerza en la fuente; no es una carga
física adicional del cilindro.

\subsection*{3.2 La inercia se cancela, pero no desaparece de la física}
\[
 \begin{aligned}
 \partial_j(u_i\tau_{ij}-v_i\sigma_{ij})
 &=-\rho\omega^2u_iv_i-u_a\delta
   +\rho\omega^2v_iu_i\\
 &=-u_a(\vct x)\delta(\vct x-\vct\xi).
 \end{aligned}
\]
\begin{idea}{Por qué puede no quedar una integral de área}
Se eligió un campo auxiliar que ya satisface el operador dinámico del
material. Por eso se cancelan los términos volumétricos de inercia.
La frecuencia sigue presente en \(U\) y \(T\); no se convirtió el problema
en uno estático.
\end{idea}

\newpage
\section*{4. Integrar y aplicar el teorema de la divergencia}
Si \(\vct\xi\) está dentro de \(\Omega\), la propiedad de la delta da
\[
 \int_\Omega u_a(\vct x)\delta(\vct x-\vct\xi)\dd\Omega
 =u_a(\vct\xi).
\]
Al integrar la divergencia:
\[
 \int_\Gamma(u_i\tau_{ij}n_j-v_i\sigma_{ij}n_j)\dd\Gamma
 =-u_a(\vct\xi).
\]
Reconociendo las tracciones \(\tau_{ij}n_j=t_i^{\mathrm{aux}}\) y
\(\sigma_{ij}n_j=t_i\), y renombrando el índice sumado:
\[
 \int_\Gamma(u_jT_{aj}-U_{aj}t_j)\dd\Gamma=-u_a(\vct\xi).
\]
Se obtiene la representación interior de Somigliana:
\begin{resultado}{El interior se representa por datos de contorno}
\[
 \boxed{u_a(\vct\xi)=
 \int_\Gamma U_{aj}(\vct x,\vct\xi)t_j(\vct x)\dd\Gamma
 -\int_\Gamma T_{aj}(\vct x,\vct\xi)u_j(\vct x)\dd\Gamma.}
\]
El segundo término lleva signo menos en la representación interior.
Al llevarlo al lado izquierdo se obtiene el signo más de la ecuación BEM.
\end{resultado}

\section*{5. Por qué la identidad todavía no es una solución}
Conocer \(U,T\) no equivale a conocer \(\vct u,\vct t\). Para este problema
se conoce \(\vct t=-p\vct n\), pero no \(\vct u\) en \(\Gamma\).
Por tanto, el segundo término sigue conteniendo funciones desconocidas.

La estrategia es imponer una ecuación de este tipo en el propio borde.
Después se aproxima la función desconocida mediante un número finito de
valores nodales. Eso corresponde a los pasos 5 y 6.

\begin{alerta}{Los productos son bilineales, no hermíticos}
En esta reciprocidad no se conjugan \(u_i\), \(v_i\), \(U\) ni \(T\).
Conjugar un estado cambiaría el operador complejo con el que se compara.
Las matrices resultantes son complejas, pero no por ello hermíticas.
\end{alerta}

\newpage
\section*{6. Llevar la fuente a la frontera}
Para construir ecuaciones nodales, se toma el límite
\(\vct\xi\to\Gamma\) desde el interior. El núcleo \(T\) es singular en
\(\vct x=\vct\xi\), de modo que no se puede sustituir sin más.

Conceptualmente se excluye una pequeña vecindad de radio \(\epsilon\),
se integra sobre el dominio recortado y se hace \(\epsilon\to0\).
La pequeña frontera creada por el recorte aporta el término libre:
\[
 \boxed{
 c_{aj}(\vct\xi)u_j(\vct\xi)+
 \VP\int_\Gamma T_{aj}u_j\dd\Gamma
 =\int_\Gamma U_{aj}t_j\dd\Gamma.}
\]
En una frontera suave,
\[
 c_{aj}=\frac12\delta_{aj}.
\]
La imagen de una ``media vecindad'' ayuda a recordar el factor \(1/2\);
su justificación precisa procede del límite singular. En esquinas el
coeficiente elástico depende de la geometría local; no es siempre
un escalar igual a un ángulo dividido por \(2\pi\).

\section*{7. Qué significa valor principal}
Una integral impropia usual integra por separado a ambos lados de la
singularidad. Un valor principal prescribe cómo se combinan los límites.
Como ejemplo escalar, no como fórmula completa del núcleo elástico:
\[
 \VP\int_{-1}^{1}\frac{1}{s}\dd s
 =\lim_{\epsilon\to0^+}
 \left(\int_{-1}^{-\epsilon}\frac{\dd s}{s}
       +\int_{\epsilon}^{1}\frac{\dd s}{s}\right)=0.
\]
Cada integral unilateral diverge, pero la combinación simétrica tiene un
límite. Cambiar arbitrariamente la exclusión puede alterar el resultado.

\begin{alerta}{Evitar un punto no define el valor principal}
No evaluar el integrando exactamente en la fuente evita una división por
cero, pero no resuelve el límite singular. La implementación combina
cuadratura adaptada con \(T=T^s+(T-T^s)\) y reconstruye la diagonal
estática mediante traslación rígida. Ese procedimiento se desarrolla
en el paso 6.
\end{alerta}

\newpage
\section*{8. Leer la ecuación como balance de influencias}
\[
 \underbrace{c\,\vct u(\vct\xi)}_{\text{término local}}
 +\underbrace{\VP\int_\Gamma T\,\vct u\,\dd\Gamma}_{\text{efecto del desplazamiento de todo el borde}}
 =\underbrace{\int_\Gamma U\,\vct t\,\dd\Gamma}_{\text{efecto de la carga conocida}}.
\]
Cada ecuación conecta un punto de observación con el contorno completo.
Eso explica por qué las matrices BEM serán densas.

\section*{9. Qué hipótesis permiten esta reducción}
\begin{itemize}
 \item Mismo material homogéneo e isótropo en el problema real y auxiliar.
 \item Linealidad, reciprocidad constitutiva y una frecuencia común.
 \item Ausencia de fuerzas de volumen en el problema real.
 Si existieran, aparecerían integrales de dominio adicionales.
 \item Uso de la normal exterior al dominio material.
 \item Tratamiento correcto de los límites singulares.
\end{itemize}
Una solución fundamental de espacio infinito puede usarse para un dominio
finito: no necesita satisfacer por sí sola la presión del cilindro.
La integral de contorno impone las condiciones reales.

\begin{resultado}{Qué gana el método y qué cuesta}
Se pasa de discretizar un área a discretizar una curva. A cambio se
requieren núcleos conocidos, integración singular y matrices globales densas.
La reducción de dimensión no garantiza menor coste en cualquier problema.
\end{resultado}

\section*{10. Comprobaciones de comprensión}
\begin{enumerate}
 \item ¿Dónde está la normal de \(T\)? En el punto \(\vct x\) de integración.
 \item ¿Por qué no hay integral de inercia de área? Porque el auxiliar
 satisface el mismo operador dinámico.
 \item ¿Dónde quedó la frecuencia? Dentro de los núcleos.
 \item ¿Ya se conoce el desplazamiento del borde? No: esta identidad es
 la ecuación que permitirá encontrarlo.
\end{enumerate}
\mapa{Paso 3: identidad continua sobre \(\Gamma\)}
{Paso 4: calcular explícitamente \(U(\vct x,\vct\xi;\omega)\) y \(T(\vct x,\vct\xi;\omega)\)}
{Pasos 5 y 6: integrar, colocar en nodos y obtener \(H\vct u=G\vct t\)}
\fuentes
\end{document}
